\chapter{Arquitectura del sistema}\label{cap:arquitectura}

\section{Vista general}

FakeNews Insight se estructura en cuatro componentes desplegables
independientes:

\begin{enumerate}[leftmargin=*,itemsep=0.2em]
  \item \textbf{Backend FastAPI} (\path{fakenews-backend/}).
  \item \textbf{Frontend React} (\path{fakenews-frontend/}).
  \item \textbf{Extensión de navegador MV3}
        (\path{browser-extension/}).
  \item \textbf{Supabase} (PostgreSQL + Auth + Storage), gestionado en
        la nube.
\end{enumerate}

% TODO: incluir diagrama de despliegue (TikZ o imagen).

\section{Patrón arquitectónico}

\subsection{Frontend: MVC adaptado}

El tutor del proyecto exige una separación estricta de
responsabilidades en el frontend. El flujo de datos es estrictamente
unidireccional:
\[
  \text{UI (componentes/páginas)}
  \;\to\;
  \text{Estado centralizado (Zustand)}
  \;\to\;
  \text{Servicios (Supabase / FastAPI)}
  \;\to\;
  \text{Base de datos}.
\]

\paragraph{Capa UI.} En \path{src/components/} y \path{src/pages/};
no hace \texttt{fetch} ni accede directamente a Supabase.

\paragraph{Capa de estado.} En \path{src/store/} con Zustand; expone
acciones (\texttt{login}, \texttt{analyze}, \texttt{verify}, \dots) que
delegan en servicios.

\paragraph{Capa de servicios.} En \path{src/services/}; único lugar que
importa \texttt{@supabase/supabase-js} o realiza \texttt{fetch} al
backend.

\subsection{Backend: separación por dominio}

El backend organiza su código por dominios:
\begin{itemize}[leftmargin=*,itemsep=0.2em]
  \item \path{main.py}: ensamblado de \emph{routers} y dependencias.
  \item \path{billing.py}: suscripciones, cuotas, integración Stripe.
  \item \path{fever/}: agente de verificación
        (\path{schemas.py}, \path{retrieval.py},
        \path{claim_extraction.py}, \path{inference.py},
        \path{aggregation.py}, \path{agent.py}).
  \item \path{tests/}: pruebas con \texttt{pytest}.
\end{itemize}

\section{Comunicación entre componentes}

\begin{itemize}[leftmargin=*,itemsep=0.2em]
  \item Frontend $\leftrightarrow$ Backend: HTTP/JSON con \emph{Bearer
        token} JWT emitido por Supabase Auth.
  \item Frontend $\leftrightarrow$ Supabase: SDK oficial.
  \item Backend $\leftrightarrow$ Supabase: PostgREST sobre PostgreSQL
        usando \emph{service role key} para operaciones administrativas.
  \item Backend $\leftrightarrow$ Stripe: API REST + recepción de
        \emph{webhooks} firmados.
  \item Extensión $\leftrightarrow$ Backend: igual que frontend
        (Bearer JWT).
\end{itemize}

\section{Decisiones transversales}

\begin{itemize}[leftmargin=*,itemsep=0.2em]
  \item \textbf{Versionado de API}: prefijo implícito \texttt{/} (no
        se versiona aún; preparado para introducir \texttt{/v1/}).
  \item \textbf{Idempotencia}: claves idempotentes en \emph{webhooks}
        de Stripe.
  \item \textbf{Configuración}: 12-factor app; variables de entorno
        documentadas en \path{.env.example}.
  \item \textbf{Errores}: respuesta uniforme JSON
        \texttt{\{detail: string\}} con códigos HTTP apropiados.
\end{itemize}
